\chapter{Spis skrótów i~symboli}

\begin{itemize}

% =========================================================
% SKRÓTY (alfabetycznie)
% =========================================================
\item[] \textbf{Skróty}

\item[16QAM] modulacja 16-QAM
\item[256QAM] modulacja 256-QAM
\item[3GPP] ang. \english{3rd Generation Partnership Project}
\item[5G] piąta generacja systemów komórkowych (ang. \english{fifth generation})
\item[64QAM] modulacja 64-QAM
\item[ALM] adaptacyjny moduł logiczny (ang. \english{adaptive logic module})
\item[ASIC] układ scalony o~specjalizowanej funkcji (ang. \english{application-specific integrated circuit})
\item[AXI4] interfejs AXI4 (ang. \english{Advanced eXtensible Interface 4})

\item[BG] graf bazowy (ang. \english{base graph})
\item[BG1] graf bazowy 1 (ang. \english{base graph 1})
\item[BG2] graf bazowy 2 (ang. \english{base graph 2})

\item[CB] blok kodowy (ang. \english{code block})
\item[CRC] suma kontrolna CRC (ang. \english{cyclic redundancy check})
\item[CRC16] CRC o~długości 16 bitów (ang. \english{CRC-16})
\item[CRC24A] CRC o~długości 24 bitów, wariant A (ang. \english{CRC-24A})
\item[CRC24B] CRC o~długości 24 bitów, wariant B (ang. \english{CRC-24B})
\item[CSR] format rzadkiej macierzy: skompresowany zapis wierszowy (ang. \english{compressed sparse row})

\item[DL-SCH] współdzielony kanał transportowy w~łączu w~dół (ang. \english{Downlink Shared Channel})
\item[DSP] przetwarzanie sygnałów cyfrowych / bloki DSP (ang. \english{digital signal processing})
\item[DUT] badany układ w~{testbench'u} (ang. \english{device under test})

\item[ETSI] Europejski Instytut Norm Telekomunikacyjnych (ang. \english{European Telecommunications Standards Institute})

\item[FPGA] programowalna macierz bramek (ang. \english{field-programmable gate array})

\item[HARQ] hybrydowy ARQ (ang. \english{hybrid automatic repeat request})
\item[HDL] język opisu sprzętu (ang. \english{hardware description language})

\item[LDPC] kody o~niskiej gęstości kontroli parzystości (ang. \english{low-density parity-check})
\item[LSB] najmniej znaczący bit (ang. \english{least significant bit})

\item[MCS] schemat modulacji i~kodowania (ang. \english{modulation and coding scheme})
\item[MSB] najbardziej znaczący bit (ang. \english{most significant bit})

\item[NR] ang. \english{New Radio}

\item[OFDM] multipleksacja z~ortogonalnym podziałem częstotliwości (ang. \english{orthogonal frequency-division multiplexing})

\item[PDSCH] fizyczny współdzielony kanał danych w~łączu w~dół (ang. \english{Physical Downlink Shared Channel})
\item[PISO] równolegle-wejściowy, szeregowo-wyjściowy (ang. \english{parallel-in serial-out})
\item[PLL] pętla synchronizacji fazowej (ang. \english{phase-locked loop})

\item[QAM] modulacja amplitudy w~kwadraturze (ang. \english{quadrature amplitude modulation})
\item[QC-LDPC] quasi-cykliczne kody LDPC (ang. \english{quasi-cyclic LDPC})
\item[QPSK] kluczowanie fazy z~przesunięciem kwadraturowym (ang. \english{quadrature phase shift keying})

\item[RAM] pamięć o~dostępie swobodnym (ang. \english{random-access memory})
\item[ROM] pamięć tylko do odczytu (ang. \english{read-only memory})
\item[RTL] poziom transferu rejestrów (ang. \english{register transfer level})
\item[RV] wersja redundancji (ang. \english{redundancy version})

\item[SDC] plik ograniczeń czasowych (ang. \english{Synopsys Design Constraints})
\item[SIPO] szeregowo-wejściowy, równolegle-wyjściowy (ang. \english{serial-in parallel-out})

\item[TB] blok transportowy (ang. \english{transport block})

\item[XOR] alternatywa wykluczająca (ang. \english{exclusive OR})

% =========================================================
% SYMBOLE (alfabetycznie)
% =========================================================
\item[] \textbf{Symbole}

\item[$A$] liczba bitów TB przed dołączeniem CRC
\item[$B$] liczba bitów po dołączeniu CRC do TB ($B=A+L$)
\item[$B'$] długość sekwencji po uwzględnieniu CRC na poziomie CB
\item[$C$] liczba bloków kodowych po segmentacji
\item[$c$] słowo kodowe (wektor wyjściowy kodera)
\item[$F$] liczba bitów wypełniających (\english{filler bits})
\item[$F_{\max}$] maksymalna częstotliwość pracy układu
\item[$g$] liczba pierwszych bloków parzystości $p_a$
\item[$H$] macierz kontroli parzystości
\item[$k_b$] parametr pośredni używany przy doborze $Z_c$
\item[$K$] długość wejściowego wektora informacji dla QC-LDPC ($K=K_b\cdot Z_c$)
\item[$K_b$] liczba bloków kolumnowych części informacyjnej (zależna od BG)
\item[$K_{CB}$] docelowa długość informacji przypadająca na pojedynczy CB po segmentacji
\item[$K_{CB\_\max}$] maksymalna długość wejściowa bloku kodowego (zależna od BG)
\item[$L$] długość CRC dołączanego do TB
\item[$L_{CB}$] długość CRC dołączanego do każdego CB
\item[$M$] długość zakodowanego strumienia po kodowaniu LDPC (przed pełnym rate matching)
\item[$N_{cykli}$] liczba cykli zegara potrzebna do przetworzenia przypadku testowego
\item[$p$] bity parzystości
\item[$p_a$] cztery pierwsze bloki parzystości
\item[$p_c$] pozostałe bloki parzystości
\item[$P_B$] liczba bloków parzystości (zależna od BG)
\item[$R_{TB}$] przepustowość/goodput dla TB
\item[$s$] część systematyczna (bity informacyjne)
\item[$T_{E2E}$] opóźnienie end-to-end (czas przetworzenia)
\item[$W$] szerokość słowa na magistrali danych (w pracy: 64 bity)
\item[$Z_c$] rząd liftingu / rozmiar podmacierzy w~QC-LDPC
\item[$Z_{c,\min}$] minimalna wartość $Z_c$ wynikająca z~długości $K_{CB}$
\item[$\alpha_{t,j}$] przesunięcie cykliczne dla składników informacji (z macierzy bazowej)
\item[$\beta_{t,j}$] przesunięcie cykliczne dla składników $p_a$ (z macierzy bazowej)
\item[$\rho$] wartość przesunięcia cyklicznego
\item[$\mathrm{rot}(\cdot,\rho)$] operator przesunięcia cyklicznego o~$\rho$

\end{itemize}
